% !TeX root = ../main.tex
\subsection{已成功计算的任务与数值结果}
\label{sec:successful-results}
为完整呈现实际计算产出，本节进一步汇总后续修复实验及历史批次中的成功计算。本文将原始输出含 ORCA 正常结束记录的调用标注为\textbf{计算成功}，并另行说明收敛、几何及性质解释的证据范围。T1--T5 与 H1 六类任务已有数值结果，H2 尚无极化率结果。这里的六类覆盖来自不同配置与多次开发尝试，不能解释为同一版本七题一次运行的成功率；它也不改变各历史批次的原始工作流完成记录。

\begin{table}[htbp]
\centering\small
\caption{V4 已获数值结果的任务覆盖。成功按实际 ORCA 调用判断，附加说明限定结论范围。}
\label{tab:successful-coverage}
\begin{tabularx}{\linewidth}{p{0.9cm}p{2.2cm}p{2.2cm}Y}
\toprule
任务 & 计算状态 & 主要批次 & 已获得结果\\
\midrule
T1 & 计算成功 & a010 & 乙醇优化能量、收敛结构与 21 个正频率\\
T2 & 计算成功 & a005、a007 & 两批共 10 个候选的单点能量及批内排序\\
T3 & 计算成功（部分点） & 第三冻结批次 & 六个约束优化计算点；角度对应待核查\\
T4 & 计算成功 & 第一冻结批次 & 丙酮五个单重态激发能及振子强度\\
T5 & 计算成功 & a026、a027 & 两种基组下各一对气相与 CPCM 水能量\\
H1 & 计算成功 & a011 & 三核总屏蔽张量及各向同性屏蔽\\
H2 & 尚无计算结果 & 已有记录 & 未获得两个取向的静态极化率\\
\bottomrule
\end{tabularx}
\end{table}

\subsubsection{T1：乙醇优化与振动频率}
批次 a010 使用气相 r$^{2}$SCAN-3c 及其内置 def2-mTZVPP 基组，电荷为 0、多重度为 1。ORCA 计算成功，优化的五项数值收敛阈值均满足，最终电子能为 $-155.002345492244\,\Eh$。27 个模式中包含投影后的 6 个平动和转动零模式、21 个正振动模式及 0 个虚频。全部正频率见表 \ref{tab:t1-frequencies}。这些结果支持该理论水平下的局部极小值，不证明全局最低构象。

\begin{table}[htbp]
\centering\small
\caption{a010 乙醇的 21 个正振动频率，单位 cm$^{-1}$，按 ORCA 打印精度列示。}
\label{tab:t1-frequencies}
\begin{tabular}{rrrrrrr}
\toprule
270.90 & 298.95 & 425.86 & 811.51 & 882.13 & 1052.13 & 1066.45\\
1134.41 & 1282.25 & 1371.13 & 1396.18 & 1418.25 & 1483.33 & 1487.57\\
1509.53 & 2969.59 & 3004.65 & 3061.83 & 3088.36 & 3105.66 & 3819.31\\
\bottomrule
\end{tabular}
\end{table}

\subsubsection{T2：正己烷有限候选的单点能量}
a007 和 a005 分别完成 6 次与 4 次气相 r$^{2}$SCAN-3c 单点计算，均为中性单重态。表 \ref{tab:t2-energies} 给出真实电子能及相对能量；每批以各自已计算子集的最低能量为参考，换算采用 $1\,\Eh=2625.4996394799$ kJ mol$^{-1}$。候选编号只在批次内部有效，相同编号不能视为同一结构。结果支持有限几何集合的单点排序，不能称为全部候选均经 DFT 优化，也不构成全局构象或自由能排序。

\begin{table}[htbp]
\centering\small
\caption{正己烷两批成功单点计算。相对能量在各批内部定义。}
\label{tab:t2-energies}
\begin{tabular}{llrr}
\toprule
批次 & 候选 & 电子能 / $\Eh$ & 相对能量 / kJ mol$^{-1}$\\
\midrule
a005 & cand\_47 & $-236.992098209035$ & 0.000000\\
a005 & cand\_26 & $-236.991452622119$ & 1.694988\\
a005 & cand\_05 & $-236.991445556634$ & 1.713539\\
a005 & cand\_02 & $-236.991444090374$ & 1.717388\\
\addlinespace
a007 & cand\_05 & $-236.992066202851$ & 0.000000\\
a007 & cand\_04 & $-236.991445556657$ & 1.629506\\
a007 & cand\_01 & $-236.991444090308$ & 1.633356\\
a007 & cand\_00 & $-236.991421448668$ & 1.692802\\
a007 & cand\_06 & $-236.991245997872$ & 2.153448\\
a007 & cand\_03 & $-236.987263819406$ & 12.608656\\
\addlinespace
\bottomrule
\end{tabular}
\end{table}

\subsubsection{T3：正丁烷约束优化的六个计算点}
第三冻结批次使用气相 r$^{2}$SCAN-3c、中性单重态，完成六次约束优化；各次均有 ORCA 正常结束和优化收敛记录，因而标为计算成功。表 \ref{tab:t3-energies} 的第一列是输入目标角度标签。生成脚本报告的实际角度与标签存在矛盾，尚未通过独立几何核查确认，不能将标签直接当作已验证二面角，也不能据此确认势垒。$+120^\circ$ 标签下的高能量尤其不能直接解释为可信物理势垒；$+180^\circ$ 点未启动。

\begin{table}[htbp]
\centering\small
\caption{成功完成的六个正丁烷计算点。角度为输入标签；相对电子能以这六个点中的最低值为零。}
\label{tab:t3-energies}
\begin{tabular}{rrr}
\toprule
输入标签 / $^\circ$ & 电子能 / $\Eh$ & 相对能量 / kJ mol$^{-1}$\\
\midrule
$-180$ & $-158.392470368263$ & 0.000000\\
$-120$ & $-158.387146187943$ & 13.978634\\
$-60$ & $-158.391526560891$ & 2.477966\\
$+0$ & $-158.383528626170$ & 23.476541\\
$+60$ & $-158.391526702130$ & 2.477595\\
$+120$ & $-158.358972460497$ & 87.948745\\
\bottomrule
\end{tabular}
\end{table}

较早批次还成功完成了一个 B3LYP/def2-TZVP 计算点。输入目标为 $-180^\circ$，收敛电子能为 $-158.397118300397\,\Eh$。独立读取终态几何得到 $179.999995786^\circ$，按角度周期性与目标一致。该点采用不同理论水平，不并入表 \ref{tab:t3-energies} 的相对能量。

\subsubsection{T4：丙酮五个单重态垂直激发}
第一冻结批次使用气相 PBE0/def2-SVP、中性单重态和固定准备几何，通过完整线性响应 TDDFT（TDA 关闭）成功计算五个单重激发态。SCF 和响应求解均有收敛记录，结果见表 \ref{tab:t4-excitations}。振子强度采用长度规范。后续产品解析未能提取这些行，不影响原始 ORCA 计算的成功判定；本表来自原始输出的只读核查。计划中的 CAM-B3LYP/def2-TZVP 没有实际执行，不能用作本表的方法标签。

\begin{table}[htbp]
\centering\small
\caption{丙酮前五个单重态垂直激发的实际 ORCA 输出。}
\label{tab:t4-excitations}
\begin{tabular}{rrrr}
\toprule
态 & 激发能 / eV & 波长 / nm & 振子强度 $f$\\
\midrule
$S_{1}$ & 4.197061 & 295.4 & 0.000007493\\
$S_{2}$ & 7.464499 & 166.1 & 0.023781213\\
$S_{3}$ & 8.406988 & 147.5 & 0.001899769\\
$S_{4}$ & 8.572306 & 144.6 & 0.060276158\\
$S_{5}$ & 8.640904 & 143.5 & 0.004121254\\
\bottomrule
\end{tabular}
\end{table}

\subsubsection{T5：固定几何下的气相与水环境电子能量}
a026 和 a027 均使用 B3LYP、中性单重态；每批气相与 CPCM 水计算使用完全相同的初始几何和原子顺序，分别完成两次 ORCA 计算。该几何不是 T1 优化后的结构。表 \ref{tab:t5-environments} 按实际基组分开报告，以 $\Delta E=E_{\mathrm{water}}-E_{\mathrm{gas}}$ 定义差值。两批结果均为负，表示指定模型与几何下水环境电子能较低；这是电子能量差，不是水合自由能。

\begin{table}[htbp]
\centering\small
\caption{两批气相与 CPCM 水环境的成功单点计算。能量单位为 $\Eh$，最后一列为 kJ mol$^{-1}$。}
\label{tab:t5-environments}
\begin{tabular}{lrrrr}
\toprule
批次 / 基组 & $E_{\mathrm{gas}}$ & $E_{\mathrm{water}}$ & $\Delta E$ & $\Delta E$\\
\midrule
a026 / def2-TZVP & $-155.002778990957$ & $-155.009740144896$ & $-0.006961153939$ & $-18.276507$\\
a027 / def2-SVP & $-154.817054999208$ & $-154.823380641772$ & $-0.006325642564$ & $-16.607972$\\
\bottomrule
\end{tabular}
\end{table}

\subsubsection{H1：HCN 的 GIAO 核磁屏蔽张量}
a011 在固定初始几何上采用气相 PBE0/def2-TZVP、中性单重态及 GIAO，ORCA 计算成功，总电子能为 $-93.345070922944\,\Eh$。按原子顺序 C$_0$、N$_1$、H$_2$，各向同性屏蔽分别为 79.051、$-34.068$ 和 29.850 ppm。完整笛卡尔总屏蔽张量如下，行列顺序为 X、Y、Z，单位 ppm，保留原始三位小数：
\begin{equation}
\bm{\sigma}_{\mathrm{C}_0}=\begin{pmatrix}
279.306 & -0.724 & -0.000 \\
-1.638 & -21.077 & 0.000 \\
-0.000 & 0.000 & -21.076
\end{pmatrix}.
\end{equation}
\begin{equation}
\bm{\sigma}_{\mathrm{N}_1}=\begin{pmatrix}
340.398 & -1.842 & -0.000 \\
-0.044 & -221.291 & 0.000 \\
0.000 & 0.000 & -221.309
\end{pmatrix}.
\end{equation}
\begin{equation}
\bm{\sigma}_{\mathrm{H}_2}=\begin{pmatrix}
39.712 & 0.055 & -0.000 \\
-0.114 & 24.919 & -0.000 \\
-0.000 & -0.000 & 24.920
\end{pmatrix}.
\end{equation}
各向同性屏蔽定义为 $\sigma_{\mathrm{iso}}=\mathrm{Tr}(\bm\sigma)/3$，与打印摘要在舍入精度内一致。上述量为绝对屏蔽，未引入参考物质，不能直接称为实验化学位移。张量非对角元按原始输出保留，不作未经说明的对称化。

\subsubsection{H2：静态极化率的结果状态}
甲醛两个取向的静态极化率张量、各向同性值及旋转一致性比较尚未得到实际 ORCA 数值，因此本项标为\textbf{尚无计算结果}。准备的几何和旋转方案不计为成功科学计算。
